\documentclass[11pt,a4paper]{article}
\usepackage[T1]{fontenc}
\usepackage[utf8]{inputenc}
\usepackage{lmodern}
\usepackage[margin=2.5cm]{geometry}
\usepackage{microtype}
\emergencystretch=1.5em
\usepackage{amsmath,amssymb,amsthm}
\usepackage{booktabs}
\usepackage{enumitem}
\usepackage[hidelinks]{hyperref}
\usepackage{url}

\newtheorem*{thmD}{Theorem D1 {\normalfont (DMV17, Theorem~1)}}
\newtheorem*{lemDtwo}{Lemma D2 {\normalfont (DMV17, Lemma~2)}}
\newtheorem*{lemDthree}{Lemma D3 {\normalfont (DMV17, Lemma~3)}}
\newtheorem*{lemtwop}{Lemma $2'$ {\normalfont (replaces Lemma~D2)}}
\newtheorem*{lemthreep}{Lemma $3'$ {\normalfont (replaces Lemma~D3)}}
\newtheorem*{thmonep}{Theorem $1'$ {\normalfont (replaces Theorem~D1)}}
\newtheorem{proposition}{Proposition}
\newtheorem{corollary}{Corollary}
\theoremstyle{definition}
\newtheorem{definition}{Definition}
\newtheorem{example}{Example}
\theoremstyle{remark}
\newtheorem{remark}{Remark}

\newcommand{\KD}{\mathsf{KD}}
\newcommand{\RD}{\mathsf{RD}}
\newcommand{\IXIF}{\mathsf{IXIF}}
\newcommand{\RO}{\mathsf{RO}}
\newcommand{\Path}{\mathit{Path}}
\newcommand{\flag}{\mathit{flag}}
\newcommand{\true}{\mathsf{true}}
\newcommand{\false}{\mathsf{false}}
\newcommand{\Z}{\mathbb{Z}}
\newcommand{\E}{\mathbb{E}}
\newcommand{\outp}[1]{\overline{#1}}
\newcommand{\inp}[1]{\widehat{#1}}
\newcommand{\mufw}{\mu_{\mathrm{fw}}}
\newcommand{\mubw}{\mu_{\mathrm{bw}}}
\newcommand{\Tfw}{T_{\mathrm{fw}}}
\newcommand{\Tbw}{T_{\mathrm{bw}}}
\newcommand{\Hmin}{H_{\min}}
\newcommand{\Hcoll}{H_{\mathrm{coll}}}
\newcommand{\NR}{\mathsf{NR}}
\newcommand{\file}[1]{\texttt{#1}}

\title{The Security Proof of the Full-State Keyed Duplex, Revisited:\\
Two Gaps, a Distinguisher Exceeding the Stated Bound, and a Repair}
\hypersetup{pdftitle={The Security Proof of the Full-State Keyed Duplex, Revisited},pdfauthor={Math\'eo Grillon}}
\author{Math\'eo Grillon\\[2pt] \normalsize Altrion-systems}
\date{17 November 2026 --- preliminary version, not peer reviewed}

\begin{document}
\maketitle

\begin{abstract}
The full-state keyed duplex of Daemen, Mennink and Van Assche (ASIACRYPT 2017) is the reference model for
permutation-based authenticated encryption, and its security bound (their Theorem~1) is used as a black box in
later work. We re-examine the proof of that theorem and report three findings.
(i) In the proof of Lemma~2 (randomized duplex versus ideal function), the number $S$ of pairs of queries whose
states can collide with probability $2^{-c}$ is bounded by $\binom{L+\Omega+1}{2}$. This count ignores that a query
whose outer part is chosen by the adversary can be aimed at a \emph{multicollision} of the outer parts that the
adversary does not control but observes. We compute the collision probability exactly and show that it exceeds the
bound written in the proof.
(ii) This is not only a gap in the proof. We describe a distinguisher that makes no query to the permutation and
whose advantage exceeds the bound of Lemma~2 and the bound of Theorem~1 itself. We prove a lower bound on its
advantage and evaluate both sides in exact rational arithmetic for five parameter sets (from $b=36$ to $b=320$,
$c=256$): the proven lower bound is $1.6$ to $2.1\,\%$ above the bound of Theorem~1. A simulation on the $36$-bit
sets agrees with the proof (three independent runs, each more than $15$ standard errors above the bound of
Theorem~1). On these instances the excess is small; it sits in a term of the form $(L+\Omega)\,\nu/2^c$ that is absent
from Theorem~1, whereas the later bound of Dobraunig and Mennink (ASIACRYPT 2019) has a term of this form.
(iii) In the proof of Lemma~3, the probability that the forward multiplicity exceeds its threshold is claimed to be at
most $\nu/2^c$ in both worlds. We prove that this fails in the real world (state recovery through the public
permutation: a lower bound that is a constant, against a claim of order $2^{-c}$), and we show by simulation only that
it fails in the randomized-duplex world (state repetition).
Finally we repair the proof. The repaired theorem keeps the structure of the original one; its bound is essentially
the original bound with the two terms of Lemma~2 counted twice, plus $2(L+\Omega)\binom{M-L}{2}/2^b$, with the
multicollision terms expressed by a direct union bound; on the five instances it is about $2.3$ times the original
bound. All numerical claims are produced by short programs, listed in the appendix and available with their outputs.
\end{abstract}

\section{Introduction}

\paragraph{Context.}
The duplex construction of Bertoni et al.~\cite{BDPV11} turns a public permutation $f$ on $b = r + c$ bits into a
stateful keyed object: each call applies $f$, outputs the $r$-bit \emph{outer part} of the state and absorbs a
block. Mennink, Reyhanitabar and Viz\'ar~\cite{MRV15} analysed the \emph{full-state} variant, in which the absorbed
block covers the whole state. Daemen, Mennink and Van Assche~\cite{DMV17} (DMV17 below) generalised it further: several
keys, a re-phased interface, and a $\flag$ input that lets a call \emph{overwrite} the outer part, as decryption
does in SpongeWrap-like modes. Their Theorem~1 bounds the advantage of distinguishing this keyed duplex $\KD$ from
an ideal extendable input function $\IXIF$ in terms of six resources $(N, M, q, q_{iv}, L, \Omega)$ and of a
\emph{multicollision limit function} $\nu$. Dobraunig and Mennink~\cite{DM19} later gave a bound of the same kind, by a
different proof, in the setting of leakage resilience. Mennink~\cite{Men23} restates both bounds side by side
(his Theorems~1 and~2) and derives the security of a number of modes from them.

We studied the proof of DMV17 while deriving the security of a duplex-based authenticated encryption mode from it.
Checking it step by step for our parameters, we met two steps that we could not justify, then found that they do
not hold as written, and finally that the first one leads to a distinguisher that contradicts the statement of the
theorem.

\paragraph{Contributions.}
In this introduction, ``Theorem~1'', ``Lemma~2'' and ``Lemma~3'' are those of DMV17.
\begin{enumerate}[leftmargin=*,itemsep=2pt]
\item \textbf{Gap 1: the count $S$ in Lemma~2} (Section~\ref{sec:gap1}). The proof bounds the probability of a state
collision in the ideal world using $S \le \binom{L+\Omega+1}{2}$. Mennink~\cite[Section~8.3]{Men23} already noted that this
step is \emph{lossy} when the adversary never sees the outer parts; we consider the general model, where it sees them,
and show by an exact computation that an adaptive adversary then \emph{exceeds} the resulting bound (Proposition~\ref{prop:f2}, Table~\ref{tab:f2}, and an example small enough
to be checked by hand).
\item \textbf{A distinguisher exceeding Lemma~2 and Theorem~1} (Section~\ref{sec:dist}). It uses $M_u$ ordinary
calls, observes a collision among their outer parts, and aims $\Omega$ overwriting calls at it. Proposition~\ref{prop:dist}
is a proven lower bound on its advantage. Table~\ref{tab:inst} compares it, in exact arithmetic, with the bounds of
Lemma~2 and Theorem~1: the lower bound is larger for all five parameter sets, from $b = 36$ to $b = 320$.
Table~\ref{tab:sim} reports the simulation on the $36$-bit sets.
\item \textbf{Gap 2: the forward multiplicity in Lemma~3} (Section~\ref{sec:gap2}). The claim
$\Pr(\mufw > \Tfw) \le \nu/2^c$ ``for both'' worlds fails in each of them. In the real world we prove
(Proposition~\ref{prop:f3}) that an adversary within the resources of the theorem reaches $\mufw > \Tfw$ with
probability about $1/(10v)$ for a small integer $v$ (for instance $0.02$ for $b = 320$, $c = 256$); in the randomized-duplex world we show it by simulation.
\item \textbf{A repair} (Section~\ref{sec:repair}): Lemma~$2'$ replaces Lemma~2, Lemma~$3'$
replaces Lemma~3, and Theorem~$1'$ replaces Theorem~1. The proofs reuse the parts of the original
proof that are not affected.
\end{enumerate}

\paragraph{What is and is not claimed.}
The statement of Theorem~1 of DMV17 does not hold as written: Proposition~\ref{prop:dist} and Table~\ref{tab:inst}
exhibit parameters, satisfying all its hypotheses, for which a distinguisher has a larger advantage. The violation
that we exhibit is \emph{quantitatively small}: on our five instances the proven lower bound is $1.6$ to $2.1\,\%$
above the bound, in a term of order $(L+\Omega)/2^c$. This is what our distinguisher achieves, not an upper bound on
what other distinguishers could achieve. It needs $L + \Omega \ge 1$: when $L = \Omega = 0$ the distinguisher does not
apply and our repaired Lemma~$2'$ coincides with the original Lemma~2. We do not claim any attack on a deployed
scheme. The repaired bound of Theorem~$1'$ is larger than the original one by a factor of about $2.3$ on these
instances (Section~\ref{sec:discussion}); this is a property of our proof, not an attack. The bound of
Dobraunig and Mennink~\cite{DM19}, proven in a different way and in a different setting, has a term of the kind that
Theorem~1 lacks (Section~\ref{sec:discussion}); we did not review that proof. Gap~2 is a gap in the proof only: the
terms of Theorem~1 cover the attacks we use to exhibit it.

\paragraph{Method and reproducibility.}
Every quantitative statement of this paper is either proven here or produced by a program listed in
Appendix~\ref{app:repro}. We keep three kinds of statements apart: \emph{proven} (Propositions~\ref{prop:f2}--\ref{prop:f3},
Lemmas~$2'$ and $3'$, Theorem~$1'$), \emph{computed exactly} from proven formulas (Tables~\ref{tab:f2} and~\ref{tab:inst},
rational arithmetic, deterministic), and \emph{measured} by Monte Carlo simulation (Table~\ref{tab:sim} and the experiments
of Section~\ref{sec:gap2}), given with their standard error, a control experiment, and a second independent
implementation. We quote DMV17 from its full version~\cite{DMV17}, by section number.
This work has not been peer reviewed. We searched for an earlier report of these points and found none in
\cite{DMV17,DM19,Men23}; Mennink~\cite[Section~8.3]{Men23} discusses the same sentence about $S$, but as a
\emph{lossy} (too pessimistic) step, in a setting where the adversary never learns the outer parts.

\section{Preliminaries}\label{sec:prelim}

\subsection{Notation}
Let $b = r + c$ with $r, c \ge 1$. For $x \in \Z_2^b$ we write $\outp{x} \in \Z_2^r$ for its first $r$ bits (the
\emph{outer part}) and $\inp{x} \in \Z_2^c$ for its last $c$ bits (the \emph{inner part}), so $x = \outp{x} \,\|\, \inp{x}$.
Addition in $\Z_2^n$ is bitwise. $\binom{n}{2} = n(n-1)/2$, and $\binom{n}{j} = 0$ for $j > n$.

\subsection{The three constructions}
We recall Algorithms~1--3 of DMV17. All have two interfaces, namely $\mathsf{Init}(\delta, iv, \sigma, \flag)$ and
$\mathsf{Duplexing}(\sigma, \flag)$, with $\delta$ a key index, $iv \in \Z_2^{b-k}$, $\sigma \in \Z_2^b$ and
$\flag \in \{\true, \false\}$; both return $Z \in \Z_2^r$.
\begin{itemize}[leftmargin=*,itemsep=2pt]
\item $\KD^f_K$ (keyed duplex) uses a $b$-bit permutation $f$ and an array $K$ of $u$ keys of $k \le b$ bits, drawn
from a distribution $D_K$. $\mathsf{Init}$ sets $s \leftarrow f(K[\delta] \,\|\, iv)$ and $\mathsf{Duplexing}$ sets
$s \leftarrow f(s)$. Then both do: $Z \leftarrow \outp{s}$; if $\flag = \true$ then $\sigma \leftarrow \sigma + (Z \| 0^c)$;
$s \leftarrow s + \sigma$; return $Z$.
\item $\RD^{\phi,\pi}$ (randomized duplex) is the same with $s \leftarrow \phi(\delta, iv)$ in $\mathsf{Init}$ and
$s \leftarrow \pi(s)$ in $\mathsf{Duplexing}$, for a uniformly random injection $\phi$ and an independent uniformly
random permutation $\pi$ of $\Z_2^b$.
\item $\IXIF^{\RO}$ keeps a string $\Path$ instead of a state. $\mathsf{Init}$ sets
$\Path \leftarrow \mathrm{Encode}(\delta) \| iv$. Then both interfaces do: $Z \leftarrow \RO(\Path)$; if $\flag = \true$
then $\sigma \leftarrow \sigma + (Z \| 0^c)$; $\Path \leftarrow \Path \| \sigma$; return $Z$. Here $\RO$ is a random oracle
with $r$-bit outputs.
\end{itemize}
In a call with $\flag = \false$ the outer part of the new state is $Z + \outp{\sigma}$: the adversary chooses
$\outp{\sigma}$ before seeing $Z$, and learns the outer part afterwards. In a call with $\flag = \true$ the outer part of
the new state is $\outp{\sigma}$: it is \emph{overwritten} with a value chosen by the adversary.

\subsection{Transcripts and resources}
Following DMV17 (Section~4.1), a $\Path$ is associated with every construction query in all three worlds, and a
query is summarised by a tuple $(\Path, Z, \sigma)$ in which $\sigma$ denotes the value \emph{after} the adjustment by
the flag. We call $t$ the state from which $Z$ is extracted ($\outp{t} = Z$) and $s = t + \sigma$ the state after the
injection. The resources of an adversary are:
$N$, the number of queries to $f$ or $f^{-1}$; $M$, the number of distinct tuples $(\Path, \sigma)$;
$q$, the number of distinct initialisation tuples $(\delta, iv)$; $q_{iv}$, the maximum number of them sharing one
$iv$; $\Omega$, the number of queries with $\flag = \true$; and $L$, which is $M$ minus the number of distinct paths.

\begin{definition}[DMV17, Definition~2]
Throw $M$ balls uniformly and independently into $2^r$ bins and let $\mu$ be the largest number of balls in a bin.
The multicollision limit function $\nu^M_{r,c}$ is the smallest natural number $x$ such that $\Pr(\mu > x) \le x/2^c$.
\end{definition}

\subsection{The theorem and the structure of its proof}
To avoid confusion with our own statements, we write Theorem~D1, Lemma~D2 and Lemma~D3 for Theorem~1, Lemma~2 and
Lemma~3 of DMV17, and Lemma~$2'$, Lemma~$3'$ and Theorem~$1'$ for the statements that replace them in Section~\ref{sec:repair}.
\begin{thmD}
Let $f$ be a random permutation, $\RO$ a random oracle and $K \leftarrow D_K$. For any adversary $A$ with resources as
above and with $N + M \le 0.1 \cdot 2^c$,
\begin{align*}
\Delta_A(\KD^f_K, f \,;\, \IXIF^{\RO}, f) \le{}& \frac{(L+\Omega)N}{2^c} + \frac{2\nu^{2(M-L)}_{r,c}(N+1)}{2^c}
 + \frac{\binom{L+\Omega+1}{2}}{2^c} \\
&+ \frac{(M-q-L)q}{2^b-q} + \frac{M(M-L-1)}{2^b} + \frac{(M-q-L)q}{2^{\Hmin(D_K)+\min\{c,\,b-k\}}} \\
&+ \frac{q_{iv}N}{2^{\Hmin(D_K)}} + \frac{\binom{u}{2}}{2^{\Hcoll(D_K)}} .
\end{align*}
\end{thmD}
For one key drawn uniformly ($u = 1$) we have $\Hmin(D_K) = k$ and the last term vanishes.
The proof (DMV17, Section~4.3) writes
$\Delta_A(\KD, f ; \IXIF, f) \le \Delta_B(\KD, f ; \RD, f) + \Delta_D(\RD ; \IXIF)$ and bounds the two terms by two lemmas.

\begin{lemDtwo}
For any adversary $D$ with resources as above and no access to $f$,
$\Delta_D(\RD^{\phi,\pi} ; \IXIF^{\RO}) \le \binom{L+\Omega+1}{2}/2^c + M(M-L-1)/2^b$.
\end{lemDtwo}
\begin{lemDthree}
For any adversary $B$ with resources as above, the advantage $\Delta_B(\KD^f_K, f ; \RD^{\phi,\pi}, f)$ is at most the
bound of Theorem~D1 without the terms $\binom{L+\Omega+1}{2}/2^c$ and $M(M-L-1)/2^b$.
\end{lemDthree}
Both lemmas are proven with the H-coefficient technique~\cite{Pat08,CS14}: if every good transcript $\tau$ satisfies
$\Pr(D_O = \tau)/\Pr(D_P = \tau) \ge 1 - \varepsilon$, then $\Delta(O ; P) \le \varepsilon + \Pr(D_P \in T_{\mathrm{bad}})$,
where $O$ is the real world and $P$ the ideal one.

\section{Gap 1: the count \texorpdfstring{$S$}{S} in the proof of Lemma~2}\label{sec:gap1}

\subsection{The step}
In the proof of Lemma~D2 (DMV17, Section~5), the transcript is extended with the states: for each path, the ideal
world draws $t = \RO(\Path) \| \RO'(\Path)$ with a second random oracle $\RO'$ with $c$-bit outputs, and the inner parts
are given to the adversary at the end. A transcript is bad if it contains a \emph{$t$-collision} (equal $t$ for
different paths) or an \emph{$s$-collision} (equal $s = t + \sigma$ for different $(\Path, \sigma)$). For good transcripts the
ratio of the H-coefficient technique is at least~$1$. The probability of an $s$-collision in the ideal world is then
bounded as follows (we quote):
\begin{quote}\small
``There are $\binom{M}{2}$ occasions for an $s$-collision. Denote by $S$ the size of the subset of these occasions for
which the adversary can (in the worst case) force the outer part of $s = t + \sigma$ to be a value of its choice. Note that
$S \le \binom{L+\Omega+1}{2}$. In the worst case, in these $S$ occasions the outer part of $s$ always has the same value
and $s$-collision probability is $2^{-c}$. For the $\binom{M}{2} - S$ other occasions the $s$-collision probability is
$2^{-b}$. Thus, the probability of an $s$-collision is upper bound by (using our bound on $S$):
\[ \frac{\binom{M}{2} - S}{2^b} + \frac{S}{2^c} \;\le\; \frac{\binom{M}{2} - \binom{L+\Omega+1}{2}}{2^b} + \frac{\binom{L+\Omega+1}{2}}{2^c}
 \;\le\; \frac{\binom{M}{2} - \binom{L+1}{2}}{2^b} + \frac{\binom{L+\Omega+1}{2}}{2^c}\,\text{.''} \]
\end{quote}
We call the last expression the \emph{written bound}; it is the largest of the three.

\subsection{What the count misses}
Call a query \emph{controlled} if its path is repeated or its flag is true (there are at most $L + \Omega$ of them),
and \emph{free} otherwise. The outer part of the state of a free query is $Z + \outp{\sigma}$ with $Z$ fresh: it is
uniform, the adversary cannot choose it, but the adversary \emph{sees} it. A controlled query made later can therefore
take as its outer part the value $\alpha$ that occurs most often among the free queries made so far. If $\mu$ free
queries share the outer part $\alpha$, the controlled query forms $\mu$ pairs, each of which is an $s$-collision with
probability $2^{-c}$, because the inner parts are independent and hidden. The count
$S \le \binom{L+\Omega+1}{2} = \binom{L+\Omega}{2} + (L+\Omega)$ allows one such pair per controlled query. It is correct
when the outer parts of the free queries are pairwise distinct, and too small otherwise.

\subsection{Exact computation}
Consider the following adversary $A_1$ against the ideal world of Lemma~D2, with parameters $M_u, \Omega \ge 1$.
It makes one initialisation and stays on one chain of calls, so all paths are distinct and $L = 0$.
\begin{enumerate}[leftmargin=*,itemsep=1pt]
\item $M_u$ calls with $\sigma = 0$ and $\flag = \false$ (the first one being $\mathsf{Init}$), with outputs $Z_1, \dots, Z_{M_u}$;
the states are $s_i = t_i$, with outer parts $Z_i$.
\item Let $\alpha$ be the most frequent value among $Z_1, \dots, Z_{M_u}$ (the smallest one in case of a tie). Then $\Omega$ calls
$\mathsf{Duplexing}(\alpha \| 0^c, \true)$; their states are $(\alpha, \inp{t})$.
\end{enumerate}
So $M = M_u + \Omega$. For the loads $m = (m_v)_{v \in \Z_2^r}$ of the outer values among the $M_u$ free queries, let
$m^* = \max_v m_v$ and $g(n) = \prod_{i=0}^{n-1}(1 - i/2^c)$.

\begin{proposition}\label{prop:f2}
In the ideal world of Lemma~D2, the transcript of $A_1$ contains an $s$-collision with probability
\[ p_1(r, c, M_u, \Omega) = 1 - \E\Bigl[\, g(m^* + \Omega) \prod_{v \in \Z_2^r,\; v \ne \alpha} g(m_v) \Bigr], \]
the expectation being over $M_u$ balls thrown uniformly and independently into $2^r$ bins, with $\alpha$ the
smallest bin of maximal load $m^* = m_\alpha$.
\end{proposition}
\begin{proof}
All paths are distinct, so the values $t_i$ are independent and uniform on $\Z_2^b$; the adversary's choices depend on
the outer parts only, so the inner parts are independent of the loads and of $\alpha$. Two states can be equal only if
they have the same outer part. Given the loads, the bin $\alpha$ receives $m^* + \Omega$ states and every other bin $v$
receives $m_v$ states, all with independent uniform inner parts; $n$ such inner parts are pairwise distinct with
probability $g(n)$, independently across bins.
\end{proof}
The program \file{f2\_exact.py} evaluates $p_1$ with rational numbers by summing over the partitions of $M_u$.
Table~\ref{tab:f2} gives the result for $r = 4$, $c = 16$, next to the written bound, which for $L = 0$ is
$\binom{M}{2}/2^b + \binom{\Omega+1}{2}/2^c$, and to the same computation for a \emph{control} adversary that draws
$\alpha$ uniformly at random.

\begin{table}[ht]
\centering
\caption{Exact probability of an $s$-collision in the ideal world of Lemma~D2, $r = 4$, $c = 16$
(\file{f2\_exact.py}). ``Written'' is the final bound written in the proof; ``Lemma~$2'$'' is the
$s$-collision part of the repaired bound of Section~\ref{sec:repair}, that is
$\binom{L+\Omega+1}{2}/2^c + \bigl[(L+\Omega)\binom{M-L}{2} + \binom{M}{2}\bigr]/2^b$ with $L = 0$ and $M = M_u + \Omega$.}\label{tab:f2}
\smallskip
\begin{tabular}{rrcccccc}
\toprule
$M_u$ & $\Omega$ & adversary $A_1$ & control & written & $A_1$/written & $\E[m^*]$ & Lemma~$2'$ \\
\midrule
 4 &  8 & $5.974\cdot10^{-4}$ & $4.634\cdot10^{-4}$ & $6.123\cdot10^{-4}$ & $0.976$ & $1.3486$ & $1.116\cdot10^{-3}$ \\
 8 &  8 & $7.056\cdot10^{-4}$ & $5.149\cdot10^{-4}$ & $6.638\cdot10^{-4}$ & $\mathbf{1.063}$ & $2.0629$ & $1.579\cdot10^{-3}$ \\
16 &  1 & $1.614\cdot10^{-4}$ & $1.297\cdot10^{-4}$ & $1.450\cdot10^{-4}$ & $\mathbf{1.113}$ & $3.0782$ & $2.747\cdot10^{-4}$ \\
16 &  8 & $9.171\cdot10^{-4}$ & $6.636\cdot10^{-4}$ & $8.125\cdot10^{-4}$ & $\mathbf{1.129}$ & $3.0782$ & $2.918\cdot10^{-3}$ \\
24 & 16 & $3.063\cdot10^{-3}$ & $2.458\cdot10^{-3}$ & $2.819\cdot10^{-3}$ & $\mathbf{1.086}$ & $3.9844$ & $1.472\cdot10^{-2}$ \\
32 &  4 & $8.591\cdot10^{-4}$ & $6.864\cdot10^{-4}$ & $7.534\cdot10^{-4}$ & $\mathbf{1.140}$ & $4.8304$ & $3.157\cdot10^{-3}$ \\
40 &  2 & $9.305\cdot10^{-4}$ & $8.351\cdot10^{-4}$ & $8.669\cdot10^{-4}$ & $\mathbf{1.073}$ & $5.6299$ & $2.509\cdot10^{-3}$ \\
 2 & 30 & $7.101\cdot10^{-3}$ & $6.674\cdot10^{-3}$ & $7.568\cdot10^{-3}$ & $0.938$ & $1.0625$ & $2.176\cdot10^{-2}$ \\
\bottomrule
\end{tabular}
\end{table}

The written bound is exceeded in six of the eight cases, by $6$ to $14\,\%$, and never by the control adversary: the
excess comes from the adaptive choice of $\alpha$. To first order in $2^{-c}$ (for fixed $r$, $M_u$, $\Omega$),
$p_1 \approx \bigl[\binom{\Omega}{2} + \Omega\,\E[m^*] + \binom{M_u}{2}/2^r\bigr]/2^c$, whereas the written bound is
$\bigl[\binom{\Omega}{2} + \Omega + \binom{M_u + \Omega}{2}/2^r\bigr]/2^c$. Since
$\binom{M_u+\Omega}{2} = \binom{M_u}{2} + M_u\Omega + \binom{\Omega}{2}$, the difference is
$\Omega\bigl(\E[m^*] - 1 - M_u/2^r\bigr)/2^c - \binom{\Omega}{2}/2^b$.

\begin{example}[checkable by hand]\label{ex:small}
Take $r = 1$, $M_u = 8$, $\Omega = 1$, so $M = 9$ and $b = c + 1$. The eight outer parts are fair bits; if $x$ of them are $0$,
the flagged query joins the larger side. With $\varepsilon = 2^{-c}$, Proposition~\ref{prop:f2} gives
$1 - p_1 = \sum_{x=0}^{8} \binom{8}{x} 2^{-8}\, g(\max(x, 8-x) + 1)\, g(\min(x, 8-x))$, that is
\[ p_1 = \tfrac{611}{32}\varepsilon - 154\varepsilon^2 + \tfrac{10927}{16}\varepsilon^3 - \tfrac{29029}{16}\varepsilon^4
 + \tfrac{94199}{32}\varepsilon^5 - \tfrac{45351}{16}\varepsilon^6 + 1473\varepsilon^7 - 315\varepsilon^8 . \]
The written bound is $\binom{9}{2}/2^{c+1} + 1/2^c = 19\varepsilon = \tfrac{608}{32}\varepsilon$. Since
$\tfrac{3}{32} > 154\varepsilon$ for $c \ge 11$ and the terms of degree $3$ to $8$ have a positive sum for $\varepsilon \le 2^{-11}$
(each positive term dominates the next one), $p_1$ exceeds the written bound for every $c \ge 11$. The program confirms
it exactly for $c = 11, 12, 16, 64$ (ratios $1.0010$, $1.0030$, $1.0048$, $1.0049$) and shows that the written bound is not exceeded for $c = 8$.
\end{example}

At this point the bound that the proof puts on the bad event is shown to be wrong. This alone does not contradict the
statement of Lemma~D2, since a bad transcript is only a sufficient condition for a gap between the two
worlds; note also that the written bound is not the right-hand side of Lemma~D2, which has $M(M-L-1)/2^b$ where the
written bound has $(\binom{M}{2} - \binom{L+1}{2})/2^b$. Table~\ref{tab:f2} therefore refutes a step of the proof, not the lemma. The next
section shows that the statement of the lemma fails as well.

\section{A distinguisher exceeding Lemma~2 and Theorem~1}\label{sec:dist}

\subsection{The distinguisher}
To turn an $s$-collision into a distinguishing event, the adversary must observe it. In $\KD$ and $\RD$ two equal states
produce equal outputs at the next call, and they stay equal as long as the same blocks are absorbed; in $\IXIF$ the
outputs on distinct paths are independent. The distinguisher $D = D_{M_u, \Omega, v}$, for integers
$M_u, \Omega, v \ge 1$, makes no query to $f$ and proceeds as follows. Its construction queries are numbered
$1, 2, \dots, M$ in order, with outputs $Z_1, \dots, Z_M$.
\begin{enumerate}[leftmargin=*,itemsep=1pt]
\item \emph{Phase 1.} Query $1$ is $\mathsf{Init}(0, 0^{b-k}, 0^b, \false)$; queries $2, \dots, M_u + v$ are
$\mathsf{Duplexing}(0^b, \false)$.
\item Let $\alpha$ be the most frequent value among $Z_1, \dots, Z_{M_u}$ (the smallest one in case of a tie) and
$T = \{ i \le M_u : Z_i = \alpha \}$.
\item \emph{Phase 2.} For $\ell = 1, \dots, \Omega$: query $a_\ell = M_u + v + 1 + (\ell - 1)v$ is
$\mathsf{Duplexing}(\alpha \| 0^c, \true)$, followed by $v - 1$ queries $\mathsf{Duplexing}(0^b, \false)$.
Then one last query $\mathsf{Duplexing}(0^b, \false)$.
\item $D$ outputs $1$ if there are $\ell \in \{1, \dots, \Omega\}$ and $x \in X_\ell = T \cup \{a_1, \dots, a_{\ell-1}\}$ with
$Z_{x+h} = Z_{a_\ell + h}$ for all $h = 1, \dots, v$; otherwise it outputs $0$.
\end{enumerate}
Its resources are $N = 0$, $q = q_{iv} = 1$, $M = M_u + v + v\Omega + 1$, exactly $\Omega$ flagged queries, and $L = 0$
since all queries lie on one chain and have distinct paths.

\subsection{A lower bound on its advantage}
Assume $M \le 2^{b-1}$ and $M_u \le 2^r$. Let $\gamma = 2^b/(2^b - M) \le 2$, and define
\begin{align*}
\pi_0 &= \prod_{i=1}^{M_u - 1} \frac{2^b - i\,2^c}{2^b - i}, \qquad \bar\mu = 1 + \gamma\binom{M_u}{2}\frac{1}{2^r}, \\
\rho &= \gamma\frac{\binom{M}{2}}{2^b} + \frac{\gamma}{2^c}\Bigl[\Omega(1 + \bar\mu) + \tbinom{\Omega}{2}
 + \frac{\gamma}{2^r}\bigl(v\Omega + (v-1)\tbinom{\Omega}{2}\bigr)\Bigr], \\
\varphi &= \Bigl[\tbinom{\Omega}{2} + \Omega\Bigl(1 + \tbinom{M_u}{2}\frac{1}{2^r}\Bigr)\Bigr]\frac{1}{2^{vr}} .
\end{align*}

\begin{proposition}\label{prop:dist}
Let $M_u, \Omega, v \ge 1$ with $M = M_u + v + v\Omega + 1 \le 2^{b-1}$ and $M_u \le 2^r$. Let $W$ be either $(\KD^f_K, f)$ with
$f$ a random permutation and one $k$-bit key $K$ ($k \le b$, any distribution), or $\RD^{\phi,\pi}$. Then
\[ \Delta_D(W \,;\, \IXIF^{\RO}) \;\ge\;
 \frac{\binom{\Omega+1}{2}(1 - \rho) + \Omega\,(1 - \pi_0 - \rho)}{2^c} \;-\; \frac{\Omega M}{2^b} \;-\; \varphi . \]
\end{proposition}
\begin{proof}
\emph{Ideal world.} All paths are distinct, so $Z_1, \dots, Z_M$ are independent and uniform. Fix $\ell$ and $x \in X_\ell$.
The indices $x + h$ ($1 \le h \le v$) are at most $a_\ell$, and $T$ depends on $Z_1, \dots, Z_{M_u}$ only; hence
$Z_{a_\ell+1}, \dots, Z_{a_\ell+v}$ are independent of $X_\ell$ and of the values they are compared with, and the $v$
equalities hold with probability $2^{-vr}$. By the union bound,
$\Pr[D^{\IXIF} = 1] \le \sum_\ell \E[|T| + \ell - 1]/2^{vr}$, and $\E|T| \le 1 + \binom{M_u}{2}/2^r$ because $|T| - 1$ is at most the
number of pairs $i < i' \le M_u$ with $Z_i = Z_{i'}$. This gives $\Pr[D^{\IXIF} = 1] \le \varphi$.

\emph{Real worlds: the sampling model.} Let $t_j$ be the state from which $Z_j$ is extracted and $s_j$ the state after the
injection: $s_j = t_j$ for the queries with $\sigma = 0$ and $s_{a_\ell} = (\alpha, \inp{t}_{a_\ell})$. In $\KD$, $t_1 = f(u_0)$ with
$u_0 = K \| 0^{b-k}$ and $t_{j+1} = f(s_j)$; in $\RD$, $t_1 = \phi(0, 0^{b-k})$ and $t_{j+1} = \pi(s_j)$. $D$ makes no query to $f$,
so these are the only evaluations of the primitives, and they are made in the order $j = 1, 2, \dots, M$. $D$ is
deterministic and its queries are fixed in advance, except for the value $\alpha$, which is a function of
$Z_1, \dots, Z_{M_u}$. Let $H_j = (K, t_1, \dots, t_j)$ be the history after query $j$ (without $K$ in $\RD$): it determines
$\alpha$ and $T$ once $j \ge M_u$, the inputs $u_0, s_1, \dots, s_j$, and the outputs $Z_1, \dots, Z_j$. Let $\NR_j$ be the event
that the inputs $u_0, s_1, \dots, s_{j-1}$ (without $u_0$ in $\RD$) are pairwise distinct; it is determined by $H_{j-1}$.

\emph{Fact~0.} For every $j \ge 1$ and every value $h$ of $H_{j-1}$ in $\NR_j$, the conditional law of $t_j$ given $H_{j-1} = h$ is
uniform on a set $U_j(h)$ with $2^b - (j-1) \le |U_j(h)| \le 2^b$: in $\KD$, $U_j = \Z_2^b \setminus \{t_1, \dots, t_{j-1}\}$; in $\RD$,
$U_1 = \Z_2^b$ and $U_j = \Z_2^b \setminus \{t_2, \dots, t_{j-1}\}$ for $j \ge 2$. Indeed, in $\KD$ the history fixes the values of the
uniformly random permutation $f$ on the $j - 1$ distinct points $u_0, s_1, \dots, s_{j-2}$, with images $t_1, \dots, t_{j-1}$, and
on $\NR_j$ the point $s_{j-1}$ (or $u_0$ if $j = 1$) is a new one; the image of a new point under a uniformly random
permutation, given its values on other points, is uniform on the values not yet taken, and $K$ is independent of $f$.
In $\RD$ the same holds for $\pi$ on the points $s_1, \dots, s_{j-1}$, and $t_1$ is uniform and independent of $\pi$.
All probabilities below are obtained by conditioning successively on $H_1, H_2, \dots$ and applying Fact~0; we write
$|U_j| \ge 2^b - M$, so that a set of $n$ values has conditional probability at most $n/(2^b - M) = \gamma n/2^b$.
We use four facts.

\emph{(a) A collision makes $D$ output $1$.} Suppose $s_{a_\ell} = s_x$ with $x \in X_\ell$. Then $t_{a_\ell+1} = t_{x+1}$.
Queries $x + 1, \dots, x + v - 1$ and $a_\ell + 1, \dots, a_\ell + v - 1$ all have $\sigma = 0$ and $\flag = \false$ (for $x \in T$
because phase~1 has $M_u + v$ such queries, for $x = a_{\ell'}$ by construction), so the states remain equal and
$t_{x+h} = t_{a_\ell+h}$ for $h \le v$. Hence $Z_{x+h} = Z_{a_\ell+h}$ for $h = 1, \dots, v$ and $D$ outputs $1$.

\emph{(b) Probability of a collision at $a_\ell$.} Let $C_\ell = \NR_{a_\ell} \wedge \{ s_{a_\ell} \in \{ s_x : x \in X_\ell \} \}$. The events
$C_\ell$ are pairwise disjoint, since $C_\ell$ contradicts $\NR_{a_{\ell'}}$ for $\ell' > \ell$. On $\NR_{a_\ell}$ the states $s_x$,
$x \in X_\ell$, are distinct and all have outer part $\alpha$ (for $x \in T$ because $s_x = t_x$ and $Z_x = \alpha$), so their inner
parts form a set $V$ of $|T| + \ell - 1$ values, determined by $H_{a_\ell-1}$, and $s_{a_\ell}$ is one of these states if and only if
$\inp{t}_{a_\ell} \in V$. Exactly $|V|\,2^r$ elements of $\Z_2^b$ have their inner part in $V$, and at most $a_\ell - 1 \le M$ of
them are outside $U_{a_\ell}$; since $|U_{a_\ell}| \le 2^b$, Fact~0 gives
$\Pr[C_\ell \mid H_{a_\ell - 1}] \ge \mathbf{1}_{\NR_{a_\ell}}\bigl((|T| + \ell - 1)/2^c - M/2^b\bigr)$.

\emph{(c) Probability that some input is repeated.} Let $\rho^* = \Pr[\neg\NR_M]$ and $I_j = \{u_0, s_1, \dots, s_{j-1}\}$, a set of
at most $j$ elements determined by $H_{j-1}$. The events $\NR_1 \supseteq \NR_2 \supseteq \dots$ are decreasing and
$\NR_{j+1} = \NR_j \wedge \{s_j \notin I_j\}$, so, decomposing along the first repetition,
\[ \rho^* = \sum_{j=1}^{M-1} \Pr[\NR_j \wedge s_j \in I_j] = \sum_{j=1}^{M-1} \E\bigl[\mathbf{1}_{\NR_j} \Pr[s_j \in I_j \mid H_{j-1}]\bigr] . \]
\emph{Queries with $\sigma = 0$.} Here $s_j = t_j$, and by Fact~0 the conditional probability is at most
$|I_j|/|U_j| \le \gamma j/2^b$. The sum over these $j$ is at most $\gamma\sum_{j=1}^{M-1} j/2^b = \gamma\binom{M}{2}/2^b$.

\emph{Flagged queries $j = a_\ell$.} Here $s_j = (\alpha, \inp{t}_j)$ with $\alpha$ determined by $H_{j-1}$, so $s_j \in I_j$ if and only
if $\inp{t}_j$ is the inner part of one of the $n_\ell$ elements of $I_j$ whose outer part is $\alpha$; each inner part
corresponds to $2^r$ values of $t_j$, so the conditional probability is at most $n_\ell 2^r/(2^b - M) = \gamma n_\ell/2^c$.
The elements of $I_{a_\ell}$ are $u_0$; the $s_i = t_i$ with $i \le M_u$, of which exactly $|T|$ have outer part $\alpha$;
the $\ell - 1$ earlier flagged states; and the states $s_{j'} = t_{j'}$ of the queries $j'$ with $M_u < j' < a_\ell$ and
$\sigma = 0$, of which there are $v + (\ell-1)(v-1)$. Hence $n_\ell \le 1 + |T| + (\ell - 1) + w_\ell$, where $w_\ell$ is the number of
the latter $j'$ with $\outp{t}_{j'} = \alpha$. For each such $j'$ we have $\NR_{a_\ell} \subseteq \NR_{j'}$, $\alpha$ is determined by
$H_{j'-1}$ because $j' > M_u$, and by Fact~0 $\Pr[\outp{t}_{j'} = \alpha \mid H_{j'-1}] \le 2^c/(2^b - M) = \gamma/2^r$ on $\NR_{j'}$;
so $\E[\mathbf{1}_{\NR_{a_\ell}} w_\ell] \le (v + (\ell-1)(v-1))\gamma/2^r$. Next, $|T| - 1$ is at most the number of pairs
$i < i' \le M_u$ with $Z_i = Z_{i'}$, and for each pair, by Fact~0 on $\NR_{i'}$,
$\Pr[Z_{i'} = Z_i \mid H_{i'-1}] \le 2^c/(2^b - M) = \gamma/2^r$; since $\NR_{a_\ell} \subseteq \NR_{i'}$,
$\E[\mathbf{1}_{\NR_{a_\ell}}|T|] \le 1 + \gamma\binom{M_u}{2}/2^r = \bar\mu$. Therefore
$\E[\mathbf{1}_{\NR_{a_\ell}} n_\ell] \le 1 + \bar\mu + (\ell - 1) + (v + (\ell-1)(v-1))\gamma/2^r$, and the sum over $\ell = 1, \dots, \Omega$ of
$\gamma/2^c$ times this quantity is $\frac{\gamma}{2^c}\bigl[\Omega(1 + \bar\mu) + \binom{\Omega}{2} + \frac{\gamma}{2^r}(v\Omega + (v-1)\binom{\Omega}{2})\bigr]$.
Adding the two sums, $\rho^* \le \rho$.

\emph{(d) Probability that $|T| = 1$.} $|T| = 1$ means that $Z_1, \dots, Z_{M_u}$ are pairwise distinct. For $1 \le i < M_u$, on
$\NR_{i+1}$ and given that $Z_1, \dots, Z_i$ are distinct, $t_{i+1}$ is uniform on $U_{i+1}$ by Fact~0. The elements of
$\Z_2^b$ whose outer part differs from $Z_1, \dots, Z_i$ number $2^b - i\,2^c \ge 0$ (as $M_u \le 2^r$) and all belong to
$U_{i+1}$, because the excluded values $t_1, \dots, t_i$ have outer parts $Z_1, \dots, Z_i$; and $|U_{i+1}| \ge 2^b - i$. So the
conditional probability that $Z_{i+1}$ is a new value is at most $(2^b - i\,2^c)/(2^b - i)$. Multiplying over $i$,
$\Pr[\NR_{M_u} \wedge |T| = 1] \le \pi_0$.

\emph{Conclusion.} By (a) and the disjointness in (b), $\Pr[D^W = 1] \ge \sum_\ell \Pr[C_\ell]$, and by (b)
$\Pr[C_\ell] \ge \E[\mathbf{1}_{\NR_{a_\ell}}(|T| + \ell - 1)]/2^c - M/2^b$. Since $|T| \ge 1$,
\[ \E[\mathbf{1}_{\NR_{a_\ell}}(|T| + \ell - 1)] \ge \ell\,\Pr[\NR_{a_\ell}] + \Pr[\NR_{a_\ell} \wedge |T| \ge 2]
 \ge \ell(1 - \rho) + (1 - \rho - \pi_0), \]
using $\Pr[\NR_{a_\ell}] \ge 1 - \rho^*$, (c) and (d). Summing over $\ell = 1, \dots, \Omega$ and subtracting the bound
$\varphi$ of the ideal world gives the claim.
\end{proof}

In words: the advantage is about $\bigl[\binom{\Omega+1}{2} + \Omega(1 - \pi_0)\bigr]/2^c$, where
$1 - \pi_0 \approx 1 - e^{-\binom{M_u}{2}/2^r}$ is the probability that two of the $M_u$ observed outer parts coincide.
For $N = 0$, $L = 0$, $q = u = 1$, Theorem~D1 claims the bound
\begin{equation}\label{eq:B}
 B = \frac{2\nu^{2M}_{r,c}}{2^c} + \frac{\binom{\Omega+1}{2}}{2^c} + \frac{M-1}{2^b-1} + \frac{M(M-1)}{2^b} + \frac{M-1}{2^{k+\min\{c,\,b-k\}}} ,
\end{equation}
and Lemma~D2 claims $B_2 = \binom{\Omega+1}{2}/2^c + M(M-1)/2^b$. Comparing with Proposition~\ref{prop:dist}, the
distinguisher exceeds $B$ when $\Omega(1 - \pi_0)$ is larger than $2\nu^{2M}_{r,c} + M(M-1)/2^r$ plus the small correction terms
of the proposition and the two terms of \eqref{eq:B} in $M - 1$ (the last of which is small only if $k$ is not much smaller than $r$). In the
instances below this is obtained with $M_u = 3 \cdot 2^{r/2}$ or a little less, and $48$ or $64$ flagged queries.

\subsection{Instances, in exact arithmetic}
The program \file{thm1\_bounds.py} evaluates $B$, $B_2$ and the lower bound of Proposition~\ref{prop:dist} with rational
numbers. In $B$ it replaces $\nu^{2M}_{r,c}$ by the upper bound $\bar\nu$, the smallest $x \ge 1$ with
$\binom{2M}{x+1}/2^{rx} \le x/2^c$: by the union bound over the $(x+1)$-subsets of balls this $x$ satisfies the
inequality that defines $\nu$, so $\nu \le \bar\nu$ and the value used for $B$ is at least the bound claimed by the
theorem. When $M_u > 10^5$ it replaces $\pi_0$ by the upper bound
$\bigl[(1 - M_u^2/2^b)\sum_{n < 60} x^n/n!\bigr]^{-1}$ with $x = \binom{M_u}{2}/2^r$, which is valid for $M_u^2 < 2^b$:
$\pi_0 = \prod_i (1 - i/2^r) \cdot \prod_i (1 - i/2^b)^{-1}$, where $\prod_i (1 - i/2^r) \le e^{-x} \le \bigl(\sum_{n<60} x^n/n!\bigr)^{-1}$ and
$\prod_{i < M_u} (1 - i/2^b) \ge (1 - M_u/2^b)^{M_u} \ge 1 - M_u^2/2^b$ (rows $4$ and $5$ of Table~\ref{tab:inst}; for row~$3$ the two values
are $0.011113$ and $0.011101$). It also checks the hypotheses $r < b$, $k \le b$ and $N + M \le 0.1 \cdot 2^c$ of Theorem~D1 and
$M \le 2^{b-1}$, $M_u \le 2^r$ of Proposition~\ref{prop:dist}.

\begin{table}[ht]
\centering
\caption{The proven lower bound $\mathit{LB}$ of Proposition~\ref{prop:dist} against the bound $B$ of Theorem~D1
(with $\nu \le \bar\nu$) and the bound $B_2$ of Lemma~D2, in units of $2^{-c}$ (\file{thm1\_bounds.py}). In all rows
$\mathit{LB} > B > B_2$ holds as an exact comparison of rational numbers.}\label{tab:inst}
\smallskip
\begin{tabular}{rrrrrrrrcccc}
\toprule
$b$ & $r$ & $c$ & $k$ & $M_u$ & $\Omega$ & $v$ & $\bar\nu$ & $B_2\cdot2^c$ & $B\cdot2^c$ & $\mathit{LB}\cdot2^c$ & $\mathit{LB}/B$ \\
\midrule
 36 & 16 &  20 &  16 & $640$ & 64 & 2 & 5 & $2089.06$ & $2099.08$ & $2135.09$ & $1.0172$ \\
 36 & 16 &  20 &  16 & $768$ & 48 & 2 & 5 & $1187.46$ & $1197.48$ & $1220.81$ & $1.0195$ \\
 64 & 24 &  40 &  24 & $3\cdot2^{12}$ & 64 & 3 & 5 & $2089.29$ & $2099.29$ & $2143.24$ & $1.0209$ \\
 96 & 32 &  64 &  32 & $3\cdot2^{16}$ & 64 & 3 & 6 & $2089.02$ & $2101.02$ & $2143.29$ & $1.0201$ \\
320 & 64 & 256 & 128 & $3\cdot2^{32}$ & 64 & 5 & 10 & $2089.00$ & $2109.00$ & $2143.29$ & $1.0163$ \\
\bottomrule
\end{tabular}
\end{table}

\begin{corollary}\label{cor:false}
For each parameter set of Table~\ref{tab:inst} there is an adversary, with resources satisfying the hypotheses of
Theorem~D1 and making no query to $f$, whose advantage in distinguishing $(\KD^f_K, f)$ from $(\IXIF^{\RO}, f)$ is
larger than the bound stated by Theorem~D1; and its advantage in distinguishing $\RD^{\phi,\pi}$ from $\IXIF^{\RO}$
is larger than the bound stated by Lemma~D2.
\end{corollary}

\subsection{Simulation}\label{sec:sim}
Corollary~\ref{cor:false} rests on Proposition~\ref{prop:dist} and on exact arithmetic, not on simulation. As an
independent check of the proposition, the program \file{thm1\_distinguisher.c} runs $D$ on the two $36$-bit parameter sets,
in the worlds $\KD$ (one uniform key, $iv = 0$), $\RD$ and $\IXIF$, with $f$ (or $\pi$) a lazily sampled uniform permutation.
The \emph{control} is the same distinguisher with $\alpha = Z_1$ instead of the most frequent value. Each line of
Table~\ref{tab:sim} is one run of $n = 4 \cdot 10^8$ independent trials; $\hat p$ is the fraction of trials with output $1$ and
its standard error is $\sqrt{\hat p(1-\hat p)/n}$. The advantage is $\Pr[D^W = 1] - \Pr[D^{\IXIF} = 1]$, and
$\Pr[D^{\IXIF} = 1] \le \varphi$ is proven, so $\hat p - \varphi$ estimates a lower bound on the advantage with the same standard
error; this is the quantity that we compare with the bounds $B$ (worlds $\KD$) and $B_2$ (world $\RD$) of Table~\ref{tab:inst}.

\begin{table}[ht]
\centering
\caption{Simulation of $D$ for $b = 36$, $r = 16$, $c = 20$, $k = 16$, $v = 2$; $4 \cdot 10^8$ trials per line
(\file{thm1\_distinguisher.c}; raw output in \file{results/thm1\_distinguisher.txt}, derived columns by \file{table3.py}).
In the column $\alpha$, ``max'' is the most frequent value and ``$Z_1$'' is the control. ``Bound'' is $B$ of Theorem~D1 in the world $\KD$ and $B_2$ of Lemma~D2 in the world $\RD$; $\mathit{LB}$ is the lower bound
of Proposition~\ref{prop:dist}; the last two columns are $(\hat p - \varphi - \text{bound})/\text{s.e.}$ and
$(\hat p - \varphi - \mathit{LB})/\text{s.e.}$, where s.e.\ is the standard error of $\hat p$. All probabilities are in units of $10^{-3}$.}\label{tab:sim}
\smallskip
\small
\setlength{\tabcolsep}{4pt}
\begin{tabular}{rrllrcccrr}
\toprule
$M_u$ & $\Omega$ & world & $\alpha$ & seed & $\hat p \pm \text{s.e.}$ & $\hat p - \varphi$ & bound & vs bound & vs $\mathit{LB}$ \\
\midrule
640 & 64 & $\KD$ & max & 1 & $2.04271 \pm 0.00226$ & $2.04218$ & $2.00184$ & $+17.8$ & $+2.7$ \\
640 & 64 & $\KD$ & max & 2 & $2.03726 \pm 0.00225$ & $2.03673$ & $2.00184$ & $+15.5$ & $+0.2$ \\
640 & 64 & $\RD$ & max & 7 & $2.04442 \pm 0.00226$ & $2.04389$ & $1.99228$ & $+22.8$ & $+3.4$ \\
640 & 64 & $\KD$ & $Z_1$ & 3 & $1.98502 \pm 0.00223$ & $1.98449$ & $2.00184$ & $-7.8$ & \\
768 & 48 & $\KD$ & max & 5 & $1.16832 \pm 0.00171$ & $1.16800$ & $1.14201$ & $+15.2$ & $+2.2$ \\
768 & 48 & $\KD$ & $Z_1$ & 6 & $1.12275 \pm 0.00167$ & $1.12243$ & $1.14201$ & $-11.7$ & \\
\bottomrule
\end{tabular}
\end{table}

In the three runs against $\KD$ with the most frequent value, the estimated advantage is $15.2$ to $17.8$ standard errors
above the bound of Theorem~D1; in the run against $\RD$ it is $22.8$ standard errors above the bound of Lemma~D2. The
control stays below the bound, which shows that the excess comes from aiming at the multicollision. The measurements
are also consistent with Proposition~\ref{prop:dist}: they lie above its lower bound, by $0.2$ to $3.4$ standard errors,
as they should since the lower bound is not tight. For the first parameter set $\varphi = 5.31 \cdot 10^{-7}$, and a run of
$4 \cdot 10^8$ trials in the world $\IXIF$ (seed $4$) gives $\Pr[D = 1] = (4.93 \pm 0.35) \cdot 10^{-7}$, below $\varphi$ as proven; for the second
set $\varphi = 3.24 \cdot 10^{-7}$. The two runs with seeds $1$ and $2$ differ by $1.7$ standard errors of their difference.
A simulation at this size cannot say anything about larger parameters; for those the statement rests on the proof alone.

\section{Gap 2: the forward multiplicity in the proof of Lemma~3}\label{sec:gap2}

\subsection{The step}
The proof of Lemma~D3 (DMV17, Section~6) defines the \emph{forward multiplicity} of a transcript of construction
queries as $\mufw = \max_\alpha \#\{ (\Path_j, Z_j, \sigma_j) : Z_j + \outp{\sigma}_j = \alpha \}$, the largest number of \emph{distinct}
tuples whose state after injection has the same outer part $\outp{s}_j = Z_j + \outp{\sigma}_j$ (as in Section~\ref{sec:prelim}, $\sigma_j$ is
the value after the adjustment by the flag, and a repeated identical query is not counted again). With $O = (\KD^f_K, f)$ and $P = (\RD^{\phi,\pi}, f)$, equation~(16) of DMV17 is
\begin{equation}\label{eq:16}
\begin{split}
\Delta_B(O ; P) \le{}& \bigl|\Pr[B^O \to 1 \wedge \mufw \le \Tfw] - \Pr[B^P \to 1 \wedge \mufw \le \Tfw]\bigr| \\
 &+ \max\bigl\{\Pr\nolimits_O(\mufw > \Tfw),\, \Pr\nolimits_P(\mufw > \Tfw)\bigr\},
\end{split}
\end{equation}
followed by: ``the bound we will derive on $\Pr(\mufw > \Tfw)$ in fact applies to both $O$ and $P$''. The first term
is bounded in a \emph{permutation setup} in which the adversary queries $\phi$ and $\pi$ directly, with at most $\Tfw$
queries to $\pi$ sharing an outer part; the resulting bound contains $\Tfw N/2^c$. Section~6.4 then sets
$\Tfw = L + \Omega + \nu^{2(M-L)}_{r,c}$ and states $\Pr(\mufw > \Tfw) \le \nu^{2(M-L)}_{r,c}/2^c$, with the justification that
``the states are generated from a set of size at least $2^b - M - N$ (for both $O$ and $P$)'', so that each of the at
most $M - L$ queries ``where the adversary has no control over the outer part'' falls in a given bin with probability
at most $2 \cdot 2^{-r}$.

The justification assumes that the output $Z_j$ of a query with a new path and $\flag = \false$ is unpredictable to
the adversary. This is not so in two situations, one in each world.

\subsection{Real world: state recovery through \texorpdfstring{$f$}{f}}
In $O$ the adversary has access to $f$. If it learns the full state of the construction, it can compute the next
output $Z$ in advance and choose $\outp{\sigma} = Z + \alpha$: the outer part is forced to $\alpha$ with $\flag = \false$ and a new
path, that is, by a query that the proof counts among those without control. Consider the adversary
$R = R_{G, v, F}$, for integers $G, v, F \ge 1$ and a fixed list $x_1, \dots, x_G$ of distinct elements of $\Z_2^c$:
\begin{enumerate}[leftmargin=*,itemsep=1pt]
\item $\mathsf{Init}(0, 0^{b-k}, 0^b, \false)$, then $v$ calls $\mathsf{Duplexing}(0^b, \false)$; outputs $Z_0, \dots, Z_v$ (in this
section the queries are numbered from $0$).
\item For $g = 1, \dots, G$: let $y_0 = Z_0 \| x_g$ and, for $i = 1, \dots, v$, $y_i = f(y_{i-1})$, stopping this $g$ as soon as
$\outp{y}_i \ne Z_i$. The first $g$ that passes the $v$ checks is \emph{accepted}, with $\mathit{cur} = y_v$.
\item If some $g$ was accepted: $F$ times, compute $p = f(\mathit{cur})$, call
$\mathsf{Duplexing}((\outp{p} + \alpha) \| 0^c, \false)$ for a fixed $\alpha$, and set $\mathit{cur} = (\alpha, \inp{p})$. Otherwise $F$
calls $\mathsf{Duplexing}(0^b, \false)$.
\end{enumerate}
Its resources are $q = 1$, $L = \Omega = 0$, $M = 1 + v + F$ and $N \le vG + F$.

\begin{proposition}\label{prop:f3}
Let $f$ be a random permutation and $K$ one $k$-bit key ($k \le b$) with any distribution. Let $G, v, F \ge 1$ with
$G \le 2^c$ and $2v + 1 \le 2^{b-1}$. Then
\[ \Pr\nolimits_O\bigl[\mufw \ge F\bigr] \;\ge\; \frac{G}{2^c} - G\Bigl(2^{(1-r)v} + \frac{4v^2}{2^b}\Bigr) - \frac{1}{2^r} - \frac{v(v+1)}{2^b} . \]
\end{proposition}
\begin{proof}
Write $u_0 = K \| 0^{b-k}$, $t_0 = f(u_0)$ and $t_i = f(t_{i-1})$, so that $Z_i = \outp{t}_i$ and the state after the $v$-th call is
$t_v$. Let $x^* = \inp{t}_0$. Since $t_0$ is uniform, $\Pr[x^* \in \{x_1, \dots, x_G\}] = G/2^c$. If $x_g = x^*$ then $y_i = t_i$ for all
$i$, so $g$ passes the $v$ checks. If $R$ accepts this $g$, then $\mathit{cur} = t_v$ is the state of the construction, and
by induction each computed $p$ is the value $f(s)$ that the construction computes in the next call, whose state after
injection is $p + ((\outp{p} + \alpha) \| 0^c) = (\alpha, \inp{p}) = \mathit{cur}$. The $F$ forced queries then all have outer part $\alpha$, and
$\mufw \ge F$. So $\mufw \ge F$ unless $x^*$ is not in the list or some $g$ with $x_g \ne x^*$ passes the $v$ checks.

Whether a given $g$ passes the $v$ checks is a property of $(K, f)$ alone: it does not depend on what $R$ did with the
guesses before $g$. We may therefore bound its probability for each $g$ separately, with a sampling order chosen for
this $g$. Fix $g$ and sample $f$ lazily in the order $u_0, t_0, \dots, t_{v-1}$, then $y_0, \dots, y_{v-1}$: the image of an input
not met before is uniform among the unused values, of which there are at least $2^b - 2v - 1 \ge 2^{b-1}$. Let $E_c$ be the
event that $u_0, t_0, \dots, t_{v-1}$ are not pairwise distinct; $\Pr[E_c] \le \sum_{i=0}^{v-1} (i+1)/(2^b - i) \le v(v+1)/2^b$.
Let $E_0$ be the event $\outp{t}_0 = \outp{u}_0$; $\Pr[E_0] = 2^{-r}$, and $y_0 = u_0$ for some $g$ implies $E_0$. Outside
$E_c \cup E_0$ and for $x_g \ne x^*$: the input $y_0$ differs from $u_0$ and $t_0$, and equals some $t_i$ with
$1 \le i \le v-1$ with probability at most $(v-1) \cdot 2/2^b$; each later input $y_{i-1}$ ($2 \le i \le v$) is a fresh
value and is one of the at most $2v + 1$ earlier inputs with probability at most $(2v+1) \cdot 2/2^b$. If all the
inputs $y_0, \dots, y_{v-1}$ are new, each check holds with probability at most $2^c/(2^b - 2v - 1) \le 2^{1-r}$ given
the past, so all $v$ hold with probability at most $2^{(1-r)v}$. Altogether a given $g$ with $x_g \ne x^*$ passes
with probability at most $2^{(1-r)v} + \bigl[2(v-1) + 2(v-1)(2v+1)\bigr]/2^b = 2^{(1-r)v} + 4(v^2-1)/2^b \le 2^{(1-r)v} + 4v^2/2^b$
outside $E_c \cup E_0$. The union bound over $g$ ends the proof.
\end{proof}

Take $F = \nu + 1$ with $\nu = \nu^{2M}_{r,c}$ and $M = 1 + v + F$, so that $\Tfw = \nu$ and the event of
Proposition~\ref{prop:f3} is $\mufw > \Tfw$, and take $G$ as large as $N + M \le 0.1 \cdot 2^c$ allows, about
$2^c/(10v)$. With $v > c/(r-1)$ the lower bound is about $1/(10v)$. The program \file{f3\_bound.py} evaluates it
exactly (with $\nu \le \bar\nu$): for $(b, r, c, v) = (36, 16, 20, 3)$, $(96, 32, 64, 3)$ and $(320, 64, 256, 5)$ it gives
$\Pr_O(\mufw > \Tfw) \ge 0.0333$, $0.0333$ and $0.0200$, against claimed bounds $\nu/2^c$ of at most
$1.91 \cdot 10^{-6}$, $1.63 \cdot 10^{-19}$ and $4.32 \cdot 10^{-77}$. This is a statement about a step of the proof under
the resources that the theorem allows: for $c = 256$ the adversary makes about $2^{252}$ queries to $f$, and nothing here
is a practical attack.

\paragraph{Simulation.} For small $r$ the proposition is weak, and we simulate instead (\file{f3\_kd.c}), with $b = 16$,
$r = 4$, $c = 12$, $k = 4$, $v = 2$, $F = 8$, hence $M = 11$, and $\Tfw = \nu^{22}_{4,12} = 7$ computed exactly by
\file{nu\_exact.py}. The simulated adversary is a \emph{budget-driven variant} of $R$, not $R_{G,v,F}$ with a fixed $G$:
it draws the guesses in a uniformly random order among all $2^c$ inner values, makes the second check of a guess only
if the first one holds, and starts a new guess only while the number $N'$ of queries to $f$ made so far satisfies
$N' + 2 + F \le 398 = 409 - M$, so that $N + M \le 409 < 0.1 \cdot 2^{12}$ in every run (the program stops with an error
otherwise). A guess costs one query, or two if the first check holds, so this variant tests about $365$ guesses,
where a fixed $G$ would have to be at most $195$ (and the static adversary is then less successful).
The adversary reaches $\mufw > 7$ with probability $0.0481 \pm 0.0005$ ($2 \cdot 10^5$ trials), $28$ times the claimed
$\nu/2^c = 0.0017$; a candidate passes the two checks in $78\,\%$ of the runs, but with $r = 4$ most of these are false
candidates, which end the search. The control (the same construction queries, no query to $f$) gives about $10^{-5}$. An
independent implementation of the same variant (\file{f3\_kd\_check.py}, Python with another generator and other seeds,
$2 \cdot 10^4$ trials) gives $0.0469 \pm 0.0015$ for the same parameters, and $0.0656 \pm 0.0018$ for $c = 10$, where $F = 7$,
$M = 10$, the budget is $N' + 2 + F \le 92$, $\nu^{20}_{4,10} = 6$ and the claimed bound is $0.0059$.

\subsection{Randomized-duplex world: state repetition}
This subsection reports an experiment; we do not prove a bound here.
In $P$ the adversary cannot recover the state through $f$, but the state of $\RD$ can \emph{repeat}: after an
$s$-collision, $\pi$ is evaluated on an input it has seen before, and the outputs repeat those of an earlier part of
the transcript. The adversary of \file{f3\_rd.c}, with parameters $\Omega$, $F$, $v$, numbers its queries from $1$ and does
the following.
\begin{enumerate}[leftmargin=*,itemsep=1pt]
\item Query $1$ is $\mathsf{Init}(0, 0^{b-k}, 0^b, \false)$; queries $2, \dots, \Omega + 1$ are $\mathsf{Duplexing}(0^b, \true)$, so that the
states after these $\Omega$ queries all have outer part $0$. Outputs $Z_1, \dots, Z_{\Omega+1}$.
\item After query $\Omega + 1$ it looks for the first pair $(j, i)$, in the order $j = 3, 4, \dots$ and, for each $j$,
$i = 2, \dots, j - 1$, such that $j + v - 1 \le \Omega + 1$ and $Z_{j+h} = Z_{i+h}$ for $h = 0, \dots, v - 1$. If there is one, $p = j - i$.
\item Queries $n = \Omega + 2, \dots, \Omega + 1 + F$: if a period $p$ was found, $\mathsf{Duplexing}(Z_{n-p} \| 0^c, \false)$, where $Z_{n-p}$
is the output that the adversary predicts for the query $n$ it is making; otherwise $\mathsf{Duplexing}(0^b, \false)$.
\end{enumerate}
Two of the $\Omega$ flagged states collide with probability about $\binom{\Omega}{2}/2^c$, and from then on the outputs are
periodic. If the period is the right one, each query of step~3 has a new path and no flag, and its state has outer
part $Z_n + Z_{n-p} = 0$, like the $\Omega$ flagged ones. Its resources are $q = 1$, $L = 0$, $N = 0$, $M = 1 + \Omega + F$ and
exactly $\Omega$ flagged queries. With $r = 8$, $c = 12$, $\Omega = 30$, $F = 6$, $v = 3$ we have $M = 37$,
$\nu^{74}_{8,12} = 5$ (\file{nu\_exact.py}) and $\Tfw = 35$.
The adversary obtains $\mufw > \Tfw$ with probability $0.0819 \pm 0.0001$ ($4 \cdot 10^6$ trials), $67$ times the claimed
$\nu/2^c = 0.0012$; a state is repeated in $10.1\,\%$ of the runs and a cycle is detected in $8.2\,\%$. The control, which
makes the same queries of steps~1 and~2 and then $F$ calls $\mathsf{Duplexing}(0^b, \false)$, gives $2.7 \cdot 10^{-5}$. An independent
implementation (\file{f3\_rd\_check.py}, Python, written from Algorithm~3 of DMV17, $3 \cdot 10^5$ trials) gives $0.0827 \pm 0.0005$.

The probability of the event is of the order of the $s$-collision probability $\binom{L+\Omega+1}{2}/2^c$, which the
claimed bound $\nu/2^c$ does not contain. The same strategy applies to $O$ with $N = 0$; we did not simulate it there.

\paragraph{Consequence.} The two attacks do not contradict Lemma~D3 or Theorem~D1: the terms
$\Tfw N/2^c$ and $\binom{L+\Omega+1}{2}/2^c$ of the theorem are larger than their success probabilities. But the
derivation of Lemma~D3 is incomplete: the term $\max\{\Pr_O, \Pr_P\}(\mufw > \Tfw)$ of \eqref{eq:16} is not bounded by
$\nu/2^c$, and in the real world it is not small at all.

\section{A repair}\label{sec:repair}

\subsection{Replacing Lemma D2}
\begin{lemtwop}
For any adversary $D$ with resources as in Section~\ref{sec:prelim} and no access to $f$,
\[ \Delta_D(\RD^{\phi,\pi} ; \IXIF^{\RO}) \le \Lambda := \frac{\binom{L+\Omega+1}{2}}{2^c}
 + \frac{(L+\Omega)\binom{M-L}{2} + \binom{M}{2} + \binom{M-L}{2}}{2^b} . \]
Moreover, the probability that the extended transcript is bad is at most $\Lambda$ in the ideal world and also in the
real world $\RD^{\phi,\pi}$.
\end{lemtwop}
\begin{proof}
As in DMV17 we may assume that $D$ is deterministic, and identical queries are counted once: the transcript is a
set of distinct tuples $(\Path, t, \sigma)$, $\sigma$ being the value after the adjustment by the flag. As in DMV17 we also use
that $\phi$ is an injection on the pairs $(\delta, iv)$ and that the path determines whether a query is an initialisation
(its path is $\mathrm{Encode}(\delta) \| iv$) or a duplexing call.
We keep the transcripts, the bad events and the ratio argument of DMV17 (Section~5): in the real world a good
transcript fixes $q$ input--output pairs of $\phi$ and $M - q - L$ pairs of $\pi$ with distinct inputs and distinct outputs,
which has probability $1/((2^b)_q (2^b)_{M-q-L}) \ge 2^{-(M-L)b}$, the probability in the ideal world
(here $(x)_n = x(x-1)\cdots(x-n+1)$). Hence
$\Pr(D_O = \tau) \ge \Pr(D_P = \tau)$ for every good $\tau$, which gives $\Delta_D \le \Pr(D_P \in T_{\mathrm{bad}})$ and also
$\Pr(D_O \in T_{\mathrm{bad}}) = 1 - \Pr(D_O \in T_{\mathrm{good}}) \le 1 - \Pr(D_P \in T_{\mathrm{good}}) = \Pr(D_P \in T_{\mathrm{bad}})$.
It remains to bound the probability of a bad transcript in the ideal world.

A $t$-collision has probability at most $\binom{M-L}{2}/2^b$, as in DMV17. For $s$-collisions, order the tuples by
the time of their first query. Two tuples with the same path have the same $t$ and different $\sigma$, hence different
$s$. The interaction depends on the outer parts only, so the inner parts $\RO'(\Path)$ of distinct paths are
independent, uniform, and independent of the interaction; two tuples with distinct paths therefore have equal inner
parts of $s$ with probability exactly $2^{-c}$, given the interaction. Hence the probability of an $s$-collision is
at most $\E[Y]/2^c$, where $Y$ is the number of pairs of tuples with distinct paths and equal outer parts of $s$.

Call a tuple \emph{free} if, at the query that creates it, its path is new (no earlier tuple has this path) and the
flag is false, and \emph{controlled} otherwise. A controlled tuple has a repeated path or was created by a flagged
query, so there are at most $L + \Omega$ of them; free tuples have distinct paths, so there are at most $M - L$. The outer part of a free tuple is $\RO(\Path) + \outp{\sigma}$
with $\RO(\Path)$ fresh and $\outp{\sigma}$ chosen before: given the past it is uniform. Consider a pair $i < j$.
If $j$ is free, the outer parts are equal with probability $2^{-r}$ given the past; these pairs contribute at most
$\binom{M}{2}/2^r$ to $\E[Y]$. If $j$ is controlled, the tuples $i < j$ with the same outer part are at most the
controlled tuples before $j$ plus $\mu_{\mathrm{free}}$, the largest number of free tuples sharing an outer part; these pairs
contribute at most $\binom{L+\Omega}{2} + (L+\Omega)\,\E[\mu_{\mathrm{free}}]$. Finally $\mu_{\mathrm{free}} - 1$ is at most the number of
pairs of free tuples with equal outer parts, so $\E[\mu_{\mathrm{free}}] \le 1 + \binom{M-L}{2}/2^r$. Adding up,
$\E[Y] \le \binom{L+\Omega+1}{2} + \bigl[(L+\Omega)\binom{M-L}{2} + \binom{M}{2}\bigr]/2^r$.
\end{proof}

For $L + \Omega = 0$, $\Lambda = M(M-1)/2^b$ is the bound of Lemma~D2. In general the new term is
$(L+\Omega)\binom{M-L}{2}/2^b$. The distinguisher of Section~\ref{sec:dist} shows that a term of this kind is necessary:
its advantage is about $\bigl[\binom{\Omega+1}{2} + \Omega(1 - e^{-\binom{M_u}{2}/2^r})\bigr]/2^c$. One can also write the
bound with $\E[\mu_{\mathrm{free}}] \le \nu^{M-L}_{r,c}\,(1 + (M-L)/2^c)$ (from $\Pr(\mu > \nu) \le \nu/2^c$ and $\mu \le M - L$), which gives a
term $(L+\Omega)\,\nu\,(1 + (M-L)/2^c)/2^c$, of the same form as the term $(L+\Omega)\nu/2^c$ of~\cite{DM19}.

\subsection{Replacing Lemma D3}
For an integer $m \ge 1$ let
\[ \beta_m = \binom{M-L}{m+1}\, 2^{(1-r)m} . \]
If at most $M - L$ balls are thrown into $2^r$ bins and each ball falls into any given bin with probability at most
$2^{1-r}$ given the previous ones, then some bin receives more than $m$ balls with probability at most $\beta_m$ (union
bound over the $(m+1)$-subsets of balls).

\begin{lemthreep}
Let $M \le 2^{b-1}$ and $m \ge 1$. For any adversary $B$ with resources as in Section~\ref{sec:prelim},
\begin{align*}
\Delta_B(\KD^f_K, f ; \RD^{\phi,\pi}, f) \le{}& \frac{(L+\Omega+2m)N}{2^c} + 2\beta_m + \Lambda + \frac{(M-q-L)q}{2^b-q} \\
 &+ \frac{(M-q-L)q}{2^{\Hmin(D_K)+\min\{c,\,b-k\}}} + \frac{q_{iv}N}{2^{\Hmin(D_K)}} + \frac{\binom{u}{2}}{2^{\Hcoll(D_K)}} .
\end{align*}
\end{lemthreep}
\begin{proof}
Let $O = (\KD^f_K, f)$, $P = (\RD^{\phi,\pi}, f)$, $\Tfw = L + \Omega + m$ and $\Tbw = m$.

\emph{Step 1: the permutation setup (unchanged).} DMV17 (Sections~6.2--6.3) shows $\KD^f_K = \RD^{I, E}$ for the
constructions $I^f_{\kappa,K}$ and $E^f_\kappa$ built from $f$ and an inner masking key $\kappa$, and proves that an adversary $C$
making $N$ queries to $f$, $q'$ to the initialisation oracle and $M'$ to the duplexing oracle, \emph{at most $\Tfw$ of
which share an outer part}, distinguishes $(I, E, f)$ from $(\phi, \pi, f)$ with advantage at most (their equation~(28))
\begin{align*}
 \varepsilon(\Tfw, \Tbw) ={}& \frac{(\Tfw + \Tbw)N}{2^c} + \Pr\nolimits_{(\phi,\pi,f)}(\mubw > \Tbw) + \frac{M'q'}{2^b - q'} \\
 &+ \frac{M'q'}{2^{\Hmin(D_K)+\min\{c,\,b-k\}}} + \frac{q_{iv}N}{2^{\Hmin(D_K)}} + \frac{\binom{u}{2}}{2^{\Hcoll(D_K)}},
\end{align*}
where $\mubw$ is the largest number of answers of $\phi$ and $\pi$ sharing an outer part. We use this result as it is.
We checked in the text of DMV17 that its derivation (Sections~6.2 and~6.3) holds for any thresholds $\Tfw$, $\Tbw$ and
uses neither the condition $N + M \le 0.1 \cdot 2^c$ nor Lemmas~4--6, which enter only in Section~6.4, where the thresholds
are chosen.
The correspondence between the two setups is the following. $C$ simulates $\RD$ for $B'$: it keeps the state $s$ of
each tuple, answers an initialisation by one query to the initialisation oracle, and answers a duplexing call
\emph{made from} a tuple with state $s$ by one query $s$ to the duplexing oracle; identical queries are answered from a
table. So the queries of $C$ to the duplexing oracle are exactly the states of the tuples that are \emph{followed by a
further call}; the state of a tuple from which no call is made is never submitted. Consequently, for any adversary
$B'$ against $(O ; P)$ with the resources of $B$ and such that, in every execution, at most $\Tfw$ of the tuples that are
followed by a further call share an outer part, the simulator $C$ satisfies the multiplicity constraint, makes the same
$N$ queries to $f$, $q' \le q$ queries to the initialisation oracle and $M' \le M - q - L$ to the duplexing oracle, and
$\Delta_{B'}(O ; P) \le \varepsilon(\Tfw, \Tbw)$. The answers of $\phi$ and $\pi$ to distinct queries are at most $M - L$ balls, each
falling in a given bin with probability at most $2^c/(2^b - M) \le 2^{1-r}$ given the previous ones; hence
$\Pr(\mubw > m) \le \beta_m$.

\emph{Step 2: truncation instead of the hybrid \eqref{eq:16}.} The quantity $\mufw$ is a function of the outputs, so $B$
can compute it. Let $B^+$ run $B$ and, as soon as $\mufw > \Tfw$, stop and output $1$; otherwise it outputs what $B$
outputs. Let $B^-$ do the same but output $0$ when it stops. Consider an execution of $B^\pm$. Before each query, $\mufw \le \Tfw$ holds for the tuples created so far, hence
for those followed by a further call. If the answer to a query creates a tuple that raises $\mufw$ to $\Tfw + 1$, the
adversary stops, so no call is made from that tuple. Thus at most $\Tfw$ tuples followed by a further call ever share
an outer part, the resources of $B^\pm$ are at most those of $B$, and Step~1 applies to $B^+$ and to $B^-$ with
$\varepsilon = \varepsilon(\Tfw, \Tbw)$. Moreover $B^\pm$ stops exactly when the full execution of $B$ would end with $\mufw > \Tfw$, since
$\mufw$ never decreases. Therefore
\begin{align*}
\Pr[B^O = 1] &\le \Pr[(B^+)^O = 1] \le \Pr[(B^+)^P = 1] + \varepsilon \le \Pr[B^P = 1] + \Pr\nolimits_P(\mufw > \Tfw) + \varepsilon, \\
\Pr[B^P = 1] &\le \Pr[(B^-)^P = 1] + \Pr\nolimits_P(\mufw > \Tfw) \le \Pr[(B^-)^O = 1] + \varepsilon + \Pr\nolimits_P(\mufw > \Tfw) \\
 &\le \Pr[B^O = 1] + \varepsilon + \Pr\nolimits_P(\mufw > \Tfw),
\end{align*}
that is, $\Delta_B(O ; P) \le \varepsilon(\Tfw, \Tbw) + \Pr_P(\mufw > \Tfw)$. Only the ideal-side probability is needed.

\emph{Step 3: the forward multiplicity in $P$.} In $P$ the oracle $f$ is independent of $\RD$; we fix $f$, so that $B$
becomes an adversary against $\RD$ alone, and the bounds below hold for every $f$, hence on average. Free and
controlled tuples are defined as in the proof of Lemma~$2'$; there are at most $L + \Omega$ controlled tuples, so
$\mufw > L + \Omega + m$ implies that more than $m$ free tuples share an outer part. Let $E$ be the event that the
transcript has an $s$-collision; $\Pr[E] \le \Lambda$ by the last statement of Lemma~$2'$. As long as no $s$-collision has
occurred, distinct tuples have distinct states, so the tuple from which a new path starts has a state that was not
submitted to $\pi$ before, and every free tuple receives a fresh answer of $\phi$ or $\pi$, uniform on at least
$2^b - M \ge 2^{b-1}$ values given the past; its outer part $Z + \outp{\sigma}$, with $\outp{\sigma}$ chosen before $Z$ is known, then
equals any given value with probability at most $2^c/(2^b - M) \le 2^{1-r}$. Number the free tuples in order of creation.
For a fixed set of $m + 1$ indices $i_0 < \dots < i_m$, the probability that the free tuples with these indices exist, that
no $s$-collision occurs before the creation of the last of them, and that they all have the outer part of the first
one, is at most $(2^{1-r})^m$, by conditioning successively on the past of the tuples $i_1, \dots, i_m$. There are at most
$\binom{M-L}{m+1}$ such sets. Hence
$\Pr_P(\mufw > \Tfw) \le \Pr[E] + \beta_m \le \Lambda + \beta_m$.

Adding $\varepsilon(L + \Omega + m, m)$ with $\Pr(\mubw > m) \le \beta_m$, $q' \le q$, $M' \le M - q - L$, and $\Lambda + \beta_m$
gives the claim.
\end{proof}

\subsection{The repaired theorem}
\begin{thmonep}
Let $f$ be a random permutation, $\RO$ a random oracle and $K \leftarrow D_K$. For any adversary $A$ with resources as in
Section~\ref{sec:prelim}, with $M \le 2^{b-1}$, and for any integer $m \ge 1$,
\begin{align*}
\Delta_A(\KD^f_K, f \,;\, \IXIF^{\RO}, f) \le{}& \frac{(L+\Omega+2m)N}{2^c} + 2\binom{M-L}{m+1}2^{(1-r)m}
 + \frac{2\binom{L+\Omega+1}{2}}{2^c} \\
&+ \frac{2\bigl[(L+\Omega)\binom{M-L}{2} + \binom{M}{2} + \binom{M-L}{2}\bigr]}{2^b} + \frac{(M-q-L)q}{2^b-q} \\
&+ \frac{(M-q-L)q}{2^{\Hmin(D_K)+\min\{c,\,b-k\}}} + \frac{q_{iv}N}{2^{\Hmin(D_K)}} + \frac{\binom{u}{2}}{2^{\Hcoll(D_K)}} .
\end{align*}
\end{thmonep}
\begin{proof}
As in DMV17 (Section~4.3), $\Delta_A(\KD, f ; \IXIF, f) \le \Delta_B(\KD, f ; \RD, f) + \Delta_D(\RD ; \IXIF)$, because $f$ is independent of
$\RD$ and of $\IXIF$. Apply Lemma~$3'$ and Lemma~$2'$.
\end{proof}

\begin{remark}[Comparison with Theorem~D1]
The integer $m$ plays the role of $\nu^{2(M-L)}_{r,c}$: the original terms $(L+\Omega)N/2^c + 2\nu(N+1)/2^c$ become
$(L+\Omega+2m)N/2^c + 2\beta_m$, and one may take for $m$ the smallest value with $\beta_m \le m/2^c$, in analogy with
the definition of $\nu$. The condition
$N + M \le 0.1 \cdot 2^c$ is no longer needed, because Lemmas~4--6 of DMV17 are not used. The remaining differences
are that the two terms of Lemma~D2 are counted twice (once for Lemma~$2'$, once for the forward
multiplicity), with $\binom{M}{2} + \binom{M-L}{2} = M(M-L-1) + \binom{L+1}{2}$ in place of $M(M-L-1)$, and that
$2(L+\Omega)\binom{M-L}{2}/2^b$ is added. The repaired bound is therefore larger than the original one. On the five
instances of Table~\ref{tab:inst} ($N = 0$, $L = 0$, and $m$ the smallest integer with $\beta_m \le m/2^c$, namely $m = 4, 4, 5, 6, 9$),
it equals $4761.97$, $2932.20$, $4773.33$, $4755.18$ and $4770.66$ times $2^{-c}$, that is $2.27$, $2.45$, $2.27$, $2.26$ and $2.26$
times the bound $B$ (\file{thm1\_bounds.py}). We do not claim that the factor $2$ on the terms of Lemma~D2 is necessary.
\end{remark}

\begin{remark}[Consistency checks]
The exact $s$-collision probabilities of Table~\ref{tab:f2} are below the $s$-collision part of $\Lambda$ in all eight
cases (last column). The lower bound of Proposition~\ref{prop:dist} and the measured values of Table~\ref{tab:sim} are below the
bound of Theorem~$1'$ on all instances (\file{thm1\_bounds.py}). These are checks of consistency, not proofs.
\end{remark}

\section{Discussion}\label{sec:discussion}

\paragraph{Size of the effect.}
The distinguisher exceeds Theorem~D1 by about $\Omega(1 - \pi_0)/2^c$ minus the slack
$2\nu/2^c + M(M-1)/2^b$ of the theorem, that is, by a few dozen times $2^{-c}$ on a bound of about $\Omega^2/2^{c+1}$:
$1.6$ to $2.1\,\%$ on the instances of Table~\ref{tab:inst}. In the other parameter choices that we examined we did not
obtain a larger relative excess with this distinguisher: a larger multicollision costs more in the term $M(M-1)/2^b$ than it
brings. This is an observation about one family of distinguishers, not a bound on the gap between Theorem~D1 and
the truth. What is established is, on one side, the lower bounds of Table~\ref{tab:inst}, and on the other side the upper bound of
Theorem~$1'$, which on the same instances is about $2.3$ times the bound of Theorem~D1, that is, a little more than one bit
of security in the terms concerned. For a given application the effect of replacing Theorem~D1 by Theorem~$1'$ has to be
computed from the resources of that application; we have not done this for any deployed scheme. When $L = \Omega = 0$,
Lemma~$2'$ and Lemma~D2 coincide and the only change is that the term $M(M-1)/2^b$ is counted twice.

\paragraph{Relation to the bound of Dobraunig and Mennink.}
The bound of~\cite{DM19}, as restated in \cite[Theorem~2]{Men23}, contains
$\bigl[\nu^M_{r,c}(L+\Omega) + \tfrac{\nu_{\mathrm{fix}}-1}{2}(L+\Omega)\bigr]/2^c$, where $\nu_{\mathrm{fix}} \le L + \Omega$ is the largest number of calls whose
outer part was set to one value. With $\nu_{\mathrm{fix}} = \Omega$ and $L = 0$ this is $\bigl[\Omega\,\nu^M_{r,c} + \binom{\Omega}{2}\bigr]/2^c$, to be
compared with the advantage $\approx \bigl[\binom{\Omega}{2} + \Omega\,\E|T|\bigr]/2^c$ of our distinguisher. The term $\Omega\nu/2^c$ has the form of the term
whose absence from Theorem~D1 the distinguisher exploits, and for the two $36$-bit instances of Table~\ref{tab:inst} it is larger than
$\Omega\,\E|T|/2^c$: there $\nu^M_{16,20} \ge 4$, because $\Pr(\mu > 3) > 3/2^{20}$ for $M$ balls in $2^{16}$ bins (a Bonferroni lower
bound, \file{thm1\_bounds.py}), whereas $\E|T| \le 2 + \binom{M_u}{3}/2^{2r} < 2.02$ for uniform outer parts, and the simulation
gives $\E|T| = 1.97$ and $2.01$ in $\KD$. We give this as an indication only. That bound is
proven by a different technique, for a construction with a different phasing of the calls and in a setting with
leakage; we did not review its proof, and we do not claim that it is a repair of Theorem~D1 in the model of DMV17.
Theorem~1 of~\cite{Men23} restates the bound of DMV17 for a slightly different phasing of the calls; we expect that it
inherits the issue, but we did not adapt the distinguisher to that phasing.

\paragraph{Limits of this work.}
(1) It has not been peer reviewed, and it has not yet been submitted to the authors of~\cite{DMV17}.
(2) In the repair we rely on the parts of the proof of DMV17 that we do not modify (the identity $\KD = \RD^{I,E}$ and the
bound~(28) in the permutation setup); we re-read their derivation and found nothing to object, but we do not reprove
them here.
(3) Proposition~\ref{prop:dist} is proven; the simulation of Section~\ref{sec:sim} is a check of it on one size, not a
substitute. Since the proposition contradicts a published theorem, an independent verification of its proof is
desirable; we wrote it so that each conditional probability refers to Fact~0.
(4) The failure of the claim of Section~\ref{sec:gap2} in the randomized-duplex world is shown by simulation only (two
implementations), not by a proof.
(5) Our search for earlier reports of these points was limited to~\cite{DMV17,DM19,Men23} and to web searches; we do
not claim that it is exhaustive.

\appendix
\section{Reproducibility}\label{app:repro}
The programs below and their raw outputs are in the repository\\
\url{https://github.com/Spider4Tech/sillage-brume}, directory\\
\file{prototypes/brume\_q/b1024/dmv17\_verif/papier}. The script \file{run\_all.sh} compiles and runs all of them and writes the files of
\file{results/}; each number of the paper is in one of these files, next to the command line that produced it.
Exact programs are deterministic. Monte Carlo programs print the estimate $\hat p$ and its standard error
$\sqrt{\hat p(1-\hat p)/n}$; their seeds are in the command lines of \file{run\_all.sh}. The program \file{thm1\_distinguisher.c}
uses one random stream per thread, so its output depends on the number of threads: the reference run used $12$.
The reference run was made with GCC~16.2.1 (\texttt{-O3 -march=native -fopenmp}), Python~3.15.0 (and Python~3.12.3 with
NumPy~2.5.4 for \file{f3\_kd\_check.py}), on an AMD Ryzen~5 5600G under Linux~7.2; it takes about one hour. A second complete run, made after the final revision of
the programs (\file{run\_all.sh}, then \file{f3\_kd\_check.py} separately in the environment that has NumPy), reproduced every
output of the reference run digit for digit; the files of \file{results/} are those of this second run.

\begin{center}
\small\begin{tabular}{lp{9.6cm}}
\toprule
File & Purpose \\
\midrule
\file{f2\_exact.py} & Proposition~\ref{prop:f2}, Table~\ref{tab:f2}, Example~\ref{ex:small} (exact, rational) \\
\file{thm1\_bounds.py} & $B$, $B_2$, the bound of Proposition~\ref{prop:dist} and the bound of Theorem~$1'$, Table~\ref{tab:inst} (exact) \\
\file{thm1\_distinguisher.c} & simulation of $D$ in $\KD$, $\RD$ and $\IXIF$ \\
\file{table3.py} & Table~\ref{tab:sim}, from the output of the simulation and from \file{thm1\_bounds.py} \\
\file{thm1\_distinguisher\_check.py} & independent implementation of $D$, written from Algorithm~1 of DMV17 \\
\file{nu\_exact.py} & exact $\nu^M_{r,c}$ for small parameters \\
\file{f3\_bound.py} & the bound of Proposition~\ref{prop:f3} (exact) \\
\file{f3\_kd.c}, \file{f3\_kd\_check.py} & gap 2 in the real world, two implementations \\
\file{f3\_rd.c}, \file{f3\_rd\_check.py} & gap 2 in the randomized-duplex world, two implementations \\
\bottomrule
\end{tabular}
\end{center}

\paragraph{Cross-check of the main simulation.} On the small model $r = 8$, $c = 10$, $k = 8$, $M_u = 40$, $\Omega = 12$, $v = 3$
($M = 80$), in the world $\KD$, the C program and the Python program, which share no code and use different generators,
give $\Pr[D = 1] = 0.08529 \pm 0.00020$ ($2 \cdot 10^6$ trials) and $0.08508 \pm 0.00051$ ($3 \cdot 10^5$ trials). This validates the
implementation only: on this small model the bound $B$ of Theorem~D1 is $0.1126$, above the measured value, and no
violation is claimed there.

\begin{thebibliography}{9}
\bibitem{BDPV11} G.~Bertoni, J.~Daemen, M.~Peeters, G.~Van Assche. Duplexing the Sponge: Single-Pass Authenticated
Encryption and Other Applications. SAC 2011, LNCS 7118, pp.~320--337. Springer, 2012.
\bibitem{CS14} S.~Chen, J.~P.~Steinberger. Tight Security Bounds for Key-Alternating Ciphers. EUROCRYPT 2014,
LNCS 8441, pp.~327--350. Springer, 2014.
\bibitem{DMV17} J.~Daemen, B.~Mennink, G.~Van Assche. Full-State Keyed Duplex with Built-In Multi-user Support.
ASIACRYPT 2017, Part~II, LNCS 10625, pp.~606--637. Springer, 2017. Full version: IACR Cryptology ePrint Archive, Report 2017/498.
\bibitem{DM19} C.~Dobraunig, B.~Mennink. Leakage Resilience of the Duplex Construction. ASIACRYPT 2019, Part~III,
LNCS 11923, pp.~225--255. Springer, 2019. Full version: IACR Cryptology ePrint Archive, Report 2019/225.
\bibitem{Men23} B.~Mennink. Understanding the Duplex and Its Security. IACR Transactions on Symmetric Cryptology,
2023(2), pp.~1--46. IACR Cryptology ePrint Archive, Report 2022/1340.
\bibitem{MRV15} B.~Mennink, R.~Reyhanitabar, D.~Viz\'ar. Security of Full-State Keyed Sponge and Duplex: Applications to
Authenticated Encryption. ASIACRYPT 2015, Part~II, LNCS 9453, pp.~465--489. Springer, 2015.
\bibitem{Pat08} J.~Patarin. The ``Coefficients H'' Technique. SAC 2008, LNCS 5381, pp.~328--345. Springer, 2008.
\end{thebibliography}
\end{document}
